\mwnewpart
\renewcommand\runningtitle{Section 3. Starting the Computer}
\chapter{Starting the Computer}
\sectioncontents

\section{From Power On to the Prompt}\ixs{startup}

\begin{plainwords}
A computer that has just been turned on knows nothing about DOS. Its BIOS
reads the first sector of the diskette, or of the fixed disk, and runs it.
That sector knows just enough to load \file{RETBIO.COM}, the part of
RetrOS/LT that works the computer's hardware. \file{RETBIO.COM} finds the
disks, loads the kernel, \file{RETDOS.COM}, and hands over to SYSINIT,
Microsoft's start-up code, which reads \file{CONFIG.SYS} and starts
\file{COMMAND.COM}. All of \file{RETBIO.COM} is new: IBM's was never
published.
\end{plainwords}

\begin{tabularx}{\textwidth}{@{}P{1.1in}P{1.2in}Y@{}}
\toprule
\textbf{Stage} & \textbf{Source} & \textbf{What happens}\\
\midrule
Master boot record & \file{boot/MBR.ASM} & On a fixed disk only. Moves
itself to 0000:0600, finds the active partition, and loads its startup
record at 0000:7C00.\\
Startup record & \file{boot/BOOT.ASM} & Checks the first two root
directory entries, loads \file{RETBIO.COM} at 0070:0000, a track at a
time.\\
INIT & \file{src/BIOS.ASM} & Takes the BIOS parameter block from the
startup record, sets the interrupt vectors and diskette parameters, finds
the drives, loads \file{RETDOS.COM}, and fills in SYSINIT's variables.\\
SYSINIT & Microsoft's \file{SYSINIT.ASM} & Moves itself to the top of
memory and the kernel down, calls the kernel's own initialization, reads
\file{CONFIG.SYS}, and loads \file{COMMAND.COM}.\\
\bottomrule
\end{tabularx}

\section{The Master Boot Record}\ixs{master boot record}

\typed{FDISK} writes \file{boot/MBR.ASM}'s code with the partition table
after it. It behaves as the IBM PC XT's: one active partition, all others
inactive; no active partition starts ROM BASIC through interrupt 18h; and
the three messages are IBM's.

\srclist{mbr_moved}{The master boot record, after it has moved.}

\section{The Startup Record}\ixs{startup record}

The PC's first BIOS does not tell the startup record which drive it was
loaded from, so \file{boot/BOOT.ASM} decides by the media byte of its own
BIOS parameter block: F8h, a fixed disk, means drive 80h; anything else,
drive 0. It copies the diskette parameter table to 0000:0522, where DOS 2
keeps it, with this disk's sectors per track, reads the first root
directory sector, and checks that \file{RETBIO.COM} and \file{RETDOS.COM}
are its first two entries. \file{RETBIO.COM} must be in one piece, as
\typed{FORMAT /S} and \typed{SYS} lay it down, so the startup record
reads it without the file allocation table:

\srclist{boot_load}{Loading RETBIO.COM, a track at a time.}

\engnote{The first version read one sector at a time. On a diskette each
such read waits nearly a full turn, 200 milliseconds, for the next sector
to come round, and \file{RETBIO.COM} is fourteen sectors. Reading to the
end of each track at once, and the same change for \file{RETDOS.COM}
below, took eight seconds off starting from the 180 KB diskette on the
5150 (Section~7).}

\section{RETBIO.COM's Devices}\ixs{RETBIO.COM}\ixs{device drivers}

\file{RETBIO.COM} is a chain of device drivers in the form the kit's
\file{DEVDRIV} document gives. \code{CON}, \code{AUX}, \code{PRN} and
\code{CLOCK\$} must be the first four, in that order:

\srclist{bios_devs}{The device headers.}

\code{CON} has the special attribute (bit 4), so the kernel writes to it
through interrupt 29h, which \file{RETBIO.COM} sends to the ROM's
teletype routine. Extended keys come back as 0 and the scan code,
\keys{Ctrl}{Break} as \keys{Ctrl}{C} (interrupt 1Bh leaves a 3 to be
read), and \keys{Ctrl}{PrtSc} as \keys{Ctrl}{P}, as MS-DOS 4's
\file{MSCON.ASM} does them. \code{CLOCK\$} turns the ROM's tick count into
the time with the arithmetic of \file{MSCLOCK.ASM}, less its AT clock.
Each driver's interrupt routine ends with \code{RETF}: the kernel calls
drivers far.

\section{The Disks}\ixs{disks}

The disk driver has two diskette units, A and B, even on a computer with
one drive, and a unit for each fixed disk with a DOS partition.

\subsection{Has the diskette changed?} The PC's drives cannot say. A
diskette used in the last two seconds is taken as unchanged; otherwise
the kernel is told ``don't know'', and reads the file allocation table
again unless it has changes waiting to be written. On a one-drive
computer, asking for the other letter means a change, after the prompt
\shown{Insert diskette for drive B: and strike any key when ready}; the
letter in the drive is kept at 0000:0504, where PC DOS keeps it.

\srclist{media_chk}{MEDIA CHECK.}

\subsection{What kind of diskette} The first byte of the file allocation
table, which is the first byte of the sector after the startup record on
all four formats, names the format: FEh, FCh, FFh or FDh. \routine{GET\_BPB}
takes the BIOS parameter block, sectors per track and heads from tables
for each.

\subsection{Reading and writing} \routine{DISKIO} works a track at a time
through the ROM, with five tries and a reset between, and maps the ROM's
error codes to the kernel's as \file{MSBDATA.INC} of MS-DOS 4 does. The
PC's DMA controller cannot transfer across a 64 KB boundary of physical
memory, and the ROM refuses a transfer that would; \routine{DISKIO}
shortens a transfer to stop before one, and passes a sector that would
straddle it through a buffer of its own:

\srclist{diskio_dma}{Never across a 64 KB boundary.}

\section{The Fixed Disk}\ixs{fixed disk}\ixs{CALCBPB}

For each fixed disk the BIOS reports, \routine{FIXEDDISK} asks its
geometry (interrupt 13h, function 8), reads its master boot record, and
looks for a partition of type 1. It does not trust the BIOS parameter
block in the partition's startup record, which does not yet exist on a
partition \typed{FDISK} has just made; it works the layout out from the
partition's size, as DOS 2 does, and \typed{FORMAT} then writes the same:

\srclist{calcbpb}{The layout of a fixed disk partition.}

512 root directory entries; two file allocation tables; and the smallest
cluster that keeps a 12-bit table within 4,084 clusters. A 10 MB XT disk
gets 8 sectors a cluster, 4,096 bytes, as PC DOS 2.10 gives it.

\section{Loading RETDOS.COM}

\file{RETDOS.COM} is the second root directory entry, and need not be in
one piece. \routine{LOADDOS} follows its chain through the first file
allocation table, reading each run of consecutive clusters with one call
to \routine{DISKIO}, to the paragraph after the whole of
\file{RETBIO.COM}:

\srclist{loaddos_run}{RETDOS.COM, a run of clusters at a time.}

\section{Handing Over to SYSINIT}\ixs{SYSINIT}

SYSINIT needs to know where the kernel is now, where it is to go, where
the device chain starts, how much memory there is, and which drive the
system started from:

\srclist{init_sysinit}{SYSINIT's variables.}

Everything from \routine{INIT} on is given back to the memory pool:
SYSINIT moves the kernel down to the paragraph after \code{RESEND}, the
end of the resident drivers.
